% SHARD-ID: AQM-35-TN-EQUALS-1-F3-ARITHMETIC-FIELD
% SHARD-TITLE: The \(t^n=1\) Arithmetic Field over \(\mathbb F_3\)
% SHARD-SUMMARY: Factors \(t^n-1\) over \(\mathbb F_3\), including the characteristic-\(3\) repeated-factor case.
% SHARD-SUMMARY: Gives the complete finite residue-field Weyl-Heisenberg representation for \(R_n=\mathbb F_3[t]/(t^n-1)\).
% SHARD-SUMMARY: Works small examples \(n=1,2,3,4,6\) and records when nilpotents are invisible to Weyl labels.
% SHARD-KEYWORDS: finite field, polynomial quotient, t^n=1, Weyl-Heisenberg, cyclotomic, nilpotents

\section{The \texorpdfstring{\(t^n=1\)}{t to the n equals 1} Arithmetic Field over \texorpdfstring{\(\mathbb F_3\)}{F3}}
\label{sec:tn-equals-1-f3-arithmetic-field}

\subsection{Status and Source Anchors}

This shard specializes AQM-34 to
\[
  R_n=\mathbb F_3[t]/(t^n-1).
\]
The input facts are the direct polynomial derivations of AQM-29 and AQM-34:
irreducible factors of \(k[t]\) give closed points, residue fields are the
corresponding finite quotients, and the Weyl representation is the tensor
product of the one-site finite-field Weyl systems from AQM-28.  No new external
source is used in this specialization.

\subsection{Lamport Dependency Skeleton}

\[
\begin{array}{rcl}
D1&:&n\ge1,\quad R_n=\mathbb F_3[t]/(t^n-1),\quad X_n=\Spec R_n.\\
D2&:&n=3^s m,\quad s=v_3(n),\quad 3\nmid m.\\
L1&:&t^n-1=(t^m-1)^{3^s},\quad t^m-1\text{ is squarefree}.\\
L2&:&\text{Closed Weyl sites are the irreducible factors of }t^m-1.\\
L3&:&\text{The degrees of these factors are determined by }
      \operatorname{ord}_d(3)\text{ for }d\mid m.\\
D3&:&E_n=\bigoplus_{\pi\mid t^m-1}K_\pi^2,\quad
      \mathcal H_n=\bigotimes_{\pi\mid t^m-1}\ell^2(K_\pi).\\
T1&:&W_n\text{ is the complete projective unitary Weyl representation on }
      \mathcal H_n.\\
T2&:&\dim_{\mathbb C}\mathcal H_n=3^m,\quad
      \mathcal A_n\cong M_{3^m}(\mathbb C).
\end{array}
\]

\subsection{Factorization in Characteristic \texorpdfstring{\(3\)}{3} L1--L3}

Let \(n=3^s m\), with \(s=v_3(n)\) and \(3\nmid m\).  In characteristic
\(3\), Frobenius gives
\[
  t^n-1=(t^m)^{3^s}-1=(t^m-1)^{3^s}.
\]
The derivative of \(t^m-1\) is \(m t^{m-1}\), and \(m\ne0\) in
\(\mathbb F_3\).  Since \(t\) does not divide \(t^m-1\), the greatest common
divisor of \(t^m-1\) and \(m t^{m-1}\) is \(1\).  Thus \(t^m-1\) is
squarefree.  Consequently the closed points of \(X_n\) are exactly the
irreducible factors of \(t^m-1\), while each such point is thickened in
\(R_n\) to multiplicity \(3^s\).

For \(d\mid m\), \(d>1\), let
\[
  o_d=\operatorname{ord}_d(3).
\]
A primitive \(d\)-th root \(\alpha\) lies in \(\mathbb F_{3^r}\) exactly when
\(\alpha^{3^r}=\alpha\), equivalently \(d\mid 3^r-1\).  Hence the degree of
\(\alpha\) over \(\mathbb F_3\) is \(o_d\).  The Frobenius orbit of
\(\alpha\) has \(o_d\) elements, so the \(\varphi(d)\) primitive \(d\)-th
roots produce
\[
  \frac{\varphi(d)}{o_d}
\]
irreducible factors of degree \(o_d\).  For \(d=1\), the factor is \(t-1\).
The sum of all factor degrees is
\[
  \sum_{d\mid m}\varphi(d)=m.
\]

\subsection{The Complete Weyl-Heisenberg Representation T1--T2}

Let \(\mathcal P_n\) be the set of monic irreducible factors of \(t^m-1\).
For \(\pi\in\mathcal P_n\), write
\[
  K_\pi=\mathbb F_3[t]/(\pi),\qquad d_\pi=\deg\pi.
\]
The closed point is \(x_\pi=(\pi)/(t^n-1)\), and its residue field is
\(\kappa(x_\pi)=K_\pi\).

Define
\[
  E_n=\bigoplus_{\pi\in\mathcal P_n}K_\pi^2,\qquad
  \mathcal H_n=\bigotimes_{\pi\in\mathcal P_n}\ell^2(K_\pi).
\]
For \(a,b,s\in K_\pi\), set
\[
  \psi_\pi(c)=
  \exp\!\left(\frac{2\pi i}{3}\Tr_{K_\pi/\mathbb F_3}(c)\right),
\]
\[
  T_\pi(a)|s\rangle=|s+a\rangle,\qquad
  R_\pi(b)|s\rangle=\psi_\pi(bs)|s\rangle,\qquad
  W_\pi(a,b)=T_\pi(a)R_\pi(b).
\]
For \(e=((a_\pi,b_\pi))_\pi\in E_n\), put
\[
  W_n(e)=\bigotimes_{\pi\in\mathcal P_n}W_\pi(a_\pi,b_\pi).
\]

\paragraph{Theorem \(T1\).}
For \(e=((a_\pi,b_\pi))\) and \(f=((a'_\pi,b'_\pi))\),
\[
  W_n(e)W_n(f)
  =
  \left(\prod_{\pi\in\mathcal P_n}\psi_\pi(b_\pi a'_\pi)\right)
  W_n(e+f),
\]
and
\[
  W_n(e)W_n(f)
  =
  \left(\prod_{\pi\in\mathcal P_n}
  \psi_\pi(b_\pi a'_\pi-b'_\pi a_\pi)\right)
  W_n(f)W_n(e).
\]
The central extension
\[
  H_n^{\rm WH}=\mu_3\times E_n,\qquad
  (\zeta,e)(\eta,f)=(\zeta\eta c_n(e,f),e+f),
\]
with
\[
  c_n(e,f)=\prod_{\pi\in\mathcal P_n}\psi_\pi(b_\pi a'_\pi),
\]
is represented unitarily by \(\rho_n(\zeta,e)=\zeta W_n(e)\).

\emph{Proof.}
This is AQM-34 applied to \(p=t^n-1\).  The one-site identity
\[
  R_\pi(b)T_\pi(a')=\psi_\pi(ba')T_\pi(a')R_\pi(b)
\]
gives the multiplier, and tensoring over all factors gives the displayed laws.

\paragraph{Theorem \(T2\).}
\[
  \operatorname{span}_{\mathbb C}\{W_n(e):e\in E_n\}
  =
  \operatorname{End}_{\mathbb C}(\mathcal H_n),
\]
so the representation is irreducible, and
\[
  \dim_{\mathbb C}\mathcal H_n
  =\prod_{\pi\in\mathcal P_n}3^{d_\pi}
  =3^m,\qquad
  \mathcal A_n\cong M_{3^m}(\mathbb C).
\]

\emph{Proof.}
At each \(\pi\), the finite-field Weyl operators span
\(\operatorname{End}(\ell^2(K_\pi))\).  Tensor products span the tensor-product
matrix algebra.  Any invariant subspace is therefore invariant under the full
matrix algebra and is \(0\) or \(\mathcal H_n\).  The dimension formula uses
\(\sum_{\pi\in\mathcal P_n}d_\pi=\deg(t^m-1)=m\).

The coordinate ring has \(3^n\) elements, but the current residue-qudit Hilbert
space has dimension \(3^m\).  They agree exactly when \(s=0\), i.e. when
\(3\nmid n\).  If \(s>0\), the nilradical is generated by the class of
\(t^m-1\), and coordinate-ring field labels factor through
\[
  R_n/(t^m-1)\cong \mathbb F_3[t]/(t^m-1).
\]

\subsection{Open Sets and States}

Because \(X_n\) is a finite set of closed points, it is discrete.  For
\(J\subset\mathcal P_n\),
\[
  U_J=\{x_\pi:\pi\in J\}
\]
is open, and
\[
  \mathcal A(U_J)=
  \bigotimes_{\pi\in J}\operatorname{End}_{\mathbb C}(\ell^2(K_\pi)).
\]
The inclusion \(J\subset J'\) sends \(a\) to
\(a\otimes\mathbf1_{J'\setminus J}\).  States on \(\mathcal A(U_J)\) are
density matrices on \(\bigotimes_{\pi\in J}\ell^2(K_\pi)\), and restriction is
partial trace over omitted factors.

\subsection{Small Cases}

For \(n=1\), \(t^n-1=t-1\).  There is one \(\mathbb F_3\) site:
\[
  \mathcal H_1\cong\mathbb C^3,\qquad
  \mathcal A_1\cong M_3(\mathbb C).
\]

For \(n=2\),
\[
  t^2-1=(t-1)(t+1),
\]
so there are two \(\mathbb F_3\) sites and
\[
  \mathcal H_2\cong\mathbb C^3\otimes\mathbb C^3,\qquad
  \mathcal A_2\cong M_9(\mathbb C).
\]

For \(n=3\),
\[
  t^3-1=(t-1)^3.
\]
The scheme is a triple thickening of the point \(t=1\), but the residue-qudit
layer is one qutrit.  If \(u=t-1\), then \(u^3=0\) in \(R_3\), and \(u\) has
zero residue at the unique Weyl site.

For \(n=4\),
\[
  t^4-1=(t-1)(t+1)(t^2+1),
\]
and \(t^2+1\) has no root in \(\mathbb F_3\).  Thus
\[
  \mathcal H_4\cong
  \mathbb C^3\otimes\mathbb C^3\otimes\mathbb C^9,\qquad
  \mathcal A_4\cong M_{81}(\mathbb C).
\]
On the \(\mathbb F_9=\mathbb F_3[u]/(u^2+1)\) factor, Frobenius sends
\(u\mapsto u^3=-u\), so
\[
  \Tr_{\mathbb F_9/\mathbb F_3}(a+bu)=2a
\]
for \(a,b\in\mathbb F_3\).  This determines the explicit phase character on
that \(9\)-level site.

For \(n=6\), \(n=3\cdot2\), so \(m=2\) and
\[
  t^6-1=(t^2-1)^3=(t-1)^3(t+1)^3.
\]
There are still exactly two residue Weyl sites, both \(\mathbb F_3\)-sites:
\[
  \mathcal H_6\cong\mathbb C^3\otimes\mathbb C^3,\qquad
  \mathcal A_6\cong M_9(\mathbb C),
\]
although the coordinate ring has \(3^6\) elements.

\subsection{Conclusion}

The \(t^n=1\) example makes the characteristic-\(3\) issue explicit.  The
residue-qudit field depends on the squarefree part \(t^m-1\), where
\(m=n/3^{v_3(n)}\).  The full coordinate ring still records the repeated
factors of \(t^n-1\) as nilpotents, but in the residue-qudit layer those
nilpotents have zero residues and therefore no separate residue Weyl
operators.  AQM-36 adds the separate nilpotent-sensitive layer.
